% ============================================================
%  IEEE Conference Paper --- IEEEtran class
%  Project 18: Water Quality Prediction --- Phase 2 (Deep Learning)
%  Semester 5 | Applied Machine Learning & Deep Learning
%  Rishihood University, Newton School of Technology
% ============================================================
\documentclass[conference]{IEEEtran}
\IEEEoverridecommandlockouts

% ── Core Packages ─────────────────────────────────────────────
\usepackage{cite}
\usepackage{amsmath,amssymb,amsfonts}
\usepackage{algorithmic}
\usepackage{graphicx}
\usepackage{textcomp}
\usepackage{xcolor}
\usepackage{booktabs}
\usepackage{multirow}
\usepackage{array}
\usepackage{url}
\usepackage{hyperref}
\usepackage{balance}
\usepackage{microtype}
\usepackage{mathptmx}
\usepackage{listings}
\usepackage{float}

\definecolor{codegreen}{rgb}{0,0.5,0}
\definecolor{codegray}{rgb}{0.4,0.4,0.4}
\definecolor{codepurple}{rgb}{0.5,0,0.5}
\definecolor{backcolour}{rgb}{0.97,0.97,0.97}

\lstdefinestyle{ieuestyle}{
    backgroundcolor=\color{backcolour},
    commentstyle=\color{codegreen}\itshape,
    keywordstyle=\color{blue}\bfseries,
    numberstyle=\tiny\color{codegray},
    stringstyle=\color{codepurple},
    basicstyle=\ttfamily\scriptsize,
    breakatwhitespace=false,
    breaklines=true,
    captionpos=b,
    keepspaces=true,
    numbers=left,
    numbersep=4pt,
    showspaces=false,
    showstringspaces=false,
    showtabs=false,
    tabsize=2,
    frame=single,
    rulecolor=\color{gray!30}
}
\lstset{style=ieuestyle}

\hypersetup{
    colorlinks  = true,
    linkcolor   = blue,
    citecolor   = blue,
    urlcolor    = blue
}

\def\BibTeX{{\rm B\kern-.05em{\sc i\kern-.025em b}\kern-.08em
    T\kern-.1667em\lower.7ex\hbox{E}\kern-.125emX}}

% ============================================================
\begin{document}
% ============================================================

\title{Phase 2: Deep Learning for Continuous Water Quality Index Prediction ---
       Multi-Layer Perceptron, Residual Hybrid Architecture,
       and Axiomatic Explainability}

\author{
  \IEEEauthorblockN{Aman Soni\IEEEauthorrefmark{1},
    Mayank Sharma\IEEEauthorrefmark{2},
    Aaryan Gera\IEEEauthorrefmark{3},\\
    Chaitanya Kumar\IEEEauthorrefmark{4},
    Ankit Kumar Pandey\IEEEauthorrefmark{5}, and
    Pranav Sehgal\IEEEauthorrefmark{6}}
  \IEEEauthorblockA{Department of Computer Science \& Artificial Intelligence,\\
    Newton School of Technology, Rishihood University\\
    Project 18: Surface Water Quality Prediction --- Semester 5\\
    Repository: \url{https://github.com/MAYANKSHARMA01010/Semester_5_AML_DL}}
}

\maketitle

% ============================================================
\begin{abstract}
Automated assessment of freshwater ecosystem health is crucial for environmental
protection, municipal water safety, and real-time response to acute pollution events.
Following the completion of Phase~1, which established tree-based baselines
(Random Forest~\cite{breiman2001random} and XGBoost~\cite{chen2016xgboost}) on a multinational corpus of 2,819,617 complete-case
observations across five nations (England, USA, Ireland, China, Canada; 1940--2023),
this paper presents the complete investigations, architectures, and empirical findings
for \textbf{Phase~2 (Deep Learning)}. 

We deliver three core methodological contributions evaluated on an identical, strictly
quarantined held-out test split of 422,943 records:
\emph{(i)} We design, optimize, and train \texttt{WaterQualityMLP}, a 5-layer deep
Multi-Layer Perceptron ($53 \to 512 \to 256 \to 128 \to 64 \to 1$) with Batch Normalization,
Dropout ($p=0.30$), and Kaiming He initialization, ingesting a 53-dimensional standardized
and one-hot encoded continuous representation;
\emph{(ii)} We construct a novel \textbf{Residual Hybrid} architecture that couples a frozen
equal-weighted Base Tree Ensemble ($0.5\times\text{RF} + 0.5\times\text{XGB}$) with a
54-dimensional neural network (\texttt{ResidualMLPNet}) trained on leakage-free 5-fold
Out-Of-Fold (OOF) residual errors under Smooth L1 (Huber) loss;
\emph{(iii)} We conduct a thorough multi-faceted empirical audit comprising systematic
7-group Leave-One-Group-Out (LOGO) feature ablations, Monte Carlo Dropout epistemic uncertainty
quantification, geographic and aquatic habitat stratification across 12 waterbody types,
computational throughput profiling, and axiomatic Integrated Gradients feature attributions
relative to a valid modal observed reference baseline ($F(x') = 86.7233$).

Empirical benchmarking demonstrates that while standalone Random Forest achieves the lowest
median absolute error ($\text{MAE} = 0.0514$, $\text{MedAE} = 0.0042$), the Residual Hybrid
model achieves the lowest global Root Mean Squared Error across all architectures
($\mathbf{\text{RMSE} = 0.1678}$, a $31.9\%$ variance reduction over the base tree ensemble)
and compresses the worst-case prediction error by \textbf{54.9\%} (MaxAE reduced from $12.6128$
down to $\mathbf{5.6842}$). LOGO ablation reveals that raw physical sensor telemetry and logarithmic
skewness normalization are irreplaceable mathematical foundations for tree and neural paradigms,
respectively. Axiomatic Integrated Gradients confirms directional physical realism, identifying
orthophosphate ($|\text{IG}| = 5.86$), nitrogen ($3.42$), and stoichiometric $\text{N}:\text{P}$
ratios ($3.04$) as the primary negative drivers of aquatic health.
\end{abstract}

\begin{IEEEkeywords}
Deep Learning, Multi-Layer Perceptron, Residual Hybrid Architecture, CCME WQI,
Water Quality Index, Integrated Gradients, Feature Ablation, Epistemic Uncertainty.
\end{IEEEkeywords}

% ============================================================
\section{Introduction}
\label{sec:intro}

\IEEEPARstart{P}{hase}~1 of this research project established foundational
machine learning benchmarks for predicting the continuous Canadian Council of Ministers
of the Environment Water Quality Index (CCME~WQI, $y \in [0, 100]$) from multi-sensor
water quality telemetry. Using 38 domain-engineered physical, stoichiometric, and
autoregressive features, Phase~1 demonstrated that tree-based ensembles (Random Forest
and XGBoost~\cite{chen2016xgboost}) achieve exceptional empirical accuracy
($R^2 > 0.999$, $\text{MAE} = 0.0514$).

However, pure tree-based models suffer from fundamental structural limitations:
\begin{enumerate}
  \item \textbf{Discretization and Boundary Artifacts:} Decision trees recursively
        partition the input feature space into axis-aligned, orthogonal hypercubes,
        producing piecewise-constant step functions. In extreme geochemical regimes,
        this produces severe outlier predictions (Random Forest exhibits a maximum
        error $\text{MaxAE} = 12.6128$).
  \item \textbf{Lack of Continuous Manifold Interpolation:} Tree models cannot
        smoothly interpolate across multivariate geochemical gradients, such as
        continuous atmospheric gas re-aeration and temperature-dependent chemical
        equilibria.
  \item \textbf{Absence of Calibrated Epistemic Uncertainty:} Standard tree ensembles
        do not provide principled, gradient-based uncertainty estimation for
        out-of-distribution environmental sensor faults.
\end{enumerate}

To overcome these challenges, \textbf{Phase~2 (Deep Learning)} explores neural network
representations, hybrid multi-paradigm ensembles, and axiomatic model interpretability
for continuous water quality prediction.

The primary contributions of Phase~2 are organized into six systematic tasks:
\begin{itemize}
  \item \textbf{Task 1: Deep Learning Data Pipeline} --- Designing a 53-dimensional
        feature representation combining 36 standardized continuous numerical features
        and 17 one-hot categorical indicators, wrapped in memory-efficient PyTorch
        \texttt{Dataset} and \texttt{DataLoader} pipelines.
  \item \textbf{Task 2: Neural Network Architecture Design} --- Implementing
        \texttt{WaterQualityMLP}, a 5-layer deep Multi-Layer Perceptron with funnel
        topology ($53 \to 512 \to 256 \to 128 \to 64 \to 1$), Batch Normalization,
        Dropout ($p=0.30$), and Kaiming He normal initialization.
  \item \textbf{Task 3: MLP Training, Optimization \& Uncertainty} --- End-to-end
        optimization using AdamW under MSE loss, learning rate decay, validation
        checkpointing, regularization sensitivity sweeps, and Monte Carlo Dropout
        epistemic uncertainty estimation with frozen BatchNorm buffers.
  \item \textbf{Task 4: Multi-Paradigm Residual Hybrid Architecture} ---
        Formulating an additive hybrid system ($\hat{y}_{\text{hybrid}} = \hat{y}_{\text{base}} + \hat{r}$)
        where a frozen Base Tree Ensemble ($0.5\times\text{RF} + 0.5\times\text{XGB}$)
        is coupled with a 54-dimensional \texttt{ResidualMLPNet} trained on leakage-free
        5-fold Out-Of-Fold (OOF) residuals under Smooth L1 (Huber) loss.
  \item \textbf{Task 5: Comprehensive Evaluation \& Ablation Studies} ---
        Unified 6-model benchmarking on 422,943 held-out test records, systematic
        7-group Leave-One-Group-Out (LOGO) feature ablations, computational inference
        latency profiling, and subgroup stratification across 5 nations and 12 aquatic habitats.
  \item \textbf{Task 6: Axiomatic Explainability \& Error Diagnostics} ---
        Path-integral Integrated Gradients feature attributions relative to a valid modal
        observed reference baseline ($F(x') = 86.7233$), signed directional chemistry
        audits, local case study decompositions, and standardized prediction error
        diagnostics ($\text{Residual} = \hat{y} - y$).
\end{itemize}

All experiments are executed on the same complete-case master corpus of 2,819,617 records
(England: 83.74\%, USA: 8.51\%, Ireland: 7.13\%, China: 0.43\%, Canada: 0.20\%) under
strict chronological and geographic anti-quarantine protocols.

% ============================================================
\section{Theoretical Background \& Related Work}
\label{sec:related}

\subsection{Phase 1 Baseline Summary}
Phase~1 trained tree-based regressors on 38 tabular features. Random Forest (100 estimators,
$\text{max\_depth}=20$) achieved $\text{MAE} = 0.0514$, $\text{RMSE} = 0.2016$, and
$R^2 = 0.999870$. XGBoost (300 estimators, $\eta = 0.05$, $\text{max\_depth}=8$) achieved
$\text{MAE} = 0.1614$, $\text{RMSE} = 0.4364$, and $R^2 = 0.999389$ with high throughput
($4.6\text{ ms}$ latency). TreeSHAP attribution~\cite{lundberg2017unified,lundberg2020from} identified orthophosphate as the dominant
predictive driver ($> 88\%$ split importance).

\subsection{Deep Learning for Tabular Environmental Telemetry}
While Deep Neural Networks (DNNs) excel in continuous domains such as vision and language,
their application to tabular data has presented challenges~\cite{gorishniy2021revisiting}.
Tabular environmental data is characterized by heterogeneous feature distributions,
varying physical scales, and sharp decision thresholds imposed by regulatory standards.
Decision trees partition such spaces efficiently via orthogonal hyperplanes, whereas
neural networks optimize continuous manifold approximations via gradient descent.
Recent advances demonstrate that incorporating Batch Normalization, Dropout, and
residual connections enables deep architectures to effectively learn non-linear tabular
representations.

\subsection{Deep Residual Learning}
He et al.~\cite{he2016deep} introduced the residual learning framework, reformulating
layers to learn residual functions $\mathcal{F}(\mathbf{x}) = \mathcal{H}(\mathbf{x}) - \mathbf{x}$
with reference to the layer inputs. In multi-paradigm regression, residual learning allows
a secondary deep learner to specialize exclusively on the errors committed by a primary
tree ensemble, targeting structured bias without disrupting accurate baseline predictions.

\subsection{Axiomatic Model Interpretability}
Sundararajan et al.~\cite{sundararajan2017axiomatic} introduced Integrated Gradients (IG),
an axiomatic attribution method satisfying two fundamental axioms:
\emph{(1) Completeness}, requiring that attributions sum to the difference between model
prediction and baseline; and \emph{(2) Implementation Invariance}, guaranteeing identical
attributions for functionally equivalent networks. For deep tabular networks, IG provides
theoretically sound attributions without the sampling variance of perturbation-based heuristics.

\subsection{Epistemic Uncertainty via Monte Carlo Dropout}
Gal and Ghahramani~\cite{gal2016dropout} developed Monte Carlo (MC) Dropout, establishing
that applying dropout during test-time inference approximates variational Bayesian inference
in deep Gaussian processes. Computing the empirical variance across $T$ stochastic forward
passes yields a practical proxy for epistemic model uncertainty.

% ============================================================
\section{Phase 2 Deep Learning Data Pipeline (Task 1)}
\label{sec:pipeline}

\subsection{53-Dimensional Input Representation}
Phase~2 expands the tabular feature space to accommodate neural network optimization
requirements. Unlike decision trees, which are invariant to monotonic transformations,
feed-forward neural networks require standardized numerical inputs and non-ordinal
categorical representations to ensure well-conditioned gradient flow:

\begin{itemize}
  \item \textbf{36 Continuous Numerical Features:} Standardized using scikit-learn
        \texttt{StandardScaler} fit exclusively on the 1,973,730 training records:
        \begin{equation}
          \tilde{x}_j = \frac{x_j - \hat{\mu}_j}{\hat{\sigma}_j + 10^{-8}}
        \end{equation}
        These 36 features encompass 8 base physical sensors, 7 stoichiometric ratios,
        7 log-transformed nutrient features, 6 station lag features, 1 environmental
        exceedance counter, and 7 temporal/cyclical markers.
  \item \textbf{17 One-Hot Encoded Categorical Features:} The two discrete identifiers
        from Phase~1 are expanded via One-Hot Encoding (\texttt{OneHotEncoder}):
        \begin{itemize}
          \item \texttt{country\_code} (5 sovereign nations: England, USA, Ireland, China,
                Canada) $\rightarrow 5$ binary indicator columns.
          \item \texttt{waterbody\_type\_code} (12 aquatic habitat classes: River, Lake,
                Canal, Estuarine, Effluent, Marine, Sea Water, Coastal, Drainage,
                Transitional, Bay, Sewage) $\rightarrow 12$ binary indicator columns.
        \end{itemize}
\end{itemize}

The resulting input vector is:
\begin{equation}
  \mathbf{x}_{\text{DL}} = \big[\,\tilde{\mathbf{x}}_{\text{num}} \;\Vert\; \mathbf{x}_{\text{country}}^{\text{OHE}} \;\Vert\; \mathbf{x}_{\text{waterbody}}^{\text{OHE}}\,\big] \in \mathbb{R}^{53}
\end{equation}

\subsection{PyTorch Dataset \& DataLoader Implementation}
A custom PyTorch \texttt{Dataset} class, \texttt{WaterQualityDataset}, was constructed
to manage memory-efficient tensor conversion and mini-batch streaming. The dataset stores
continuous feature matrices as single-precision floating-point tensors (\texttt{torch.float32})
and regression targets as 1D scalar tensors.

DataLoaders were configured as follows:
\begin{itemize}
  \item \textbf{Training Loader:} $\text{Batch Size} = 2,048$, pseudo-random shuffling enabled
        (\texttt{seed=42}), 4 background worker threads, pinned memory enabled.
  \item \textbf{Validation / Test Loaders:} $\text{Batch Size} = 4,096$, deterministic sequential
        sampling (no shuffling), 4 worker threads.
\end{itemize}

Zero data leakage was programmatically enforced: scalers and encoders were fit strictly on the
training split and serialized to \texttt{artifacts/} for application to validation and test sets.

% ============================================================
\section{WaterQualityMLP Architecture Design (Task 2)}
\label{sec:mlp_arch}

\subsection{Funnel Topology Specification}
The primary deep learning architecture, \texttt{WaterQualityMLP}, implements a 5-layer
fully connected feed-forward funnel topology:
\begin{equation}
  53 \xrightarrow{L_1} 512 \xrightarrow{L_2} 256 \xrightarrow{L_3} 128 \xrightarrow{L_4} 64 \xrightarrow{L_5} 1
\end{equation}

The funnel structure progressively compresses high-dimensional multi-sensor interactions
into compact, abstract latent representations before producing the scalar continuous WQI prediction.

\subsection{Layer-Wise Regularization \& Activation Block}
Each hidden layer block $l \in \{1, 2, 3, 4\}$ applies the following sequential operations:
\begin{equation}
  \mathbf{z}^{(l)} = \mathbf{W}^{(l)}\mathbf{h}^{(l-1)} + \mathbf{b}^{(l)}
\end{equation}
\begin{equation}
  \hat{\mathbf{z}}^{(l)} = \text{BatchNorm1d}\!\left(\mathbf{z}^{(l)}\right) = \gamma^{(l)} \odot \frac{\mathbf{z}^{(l)} - \mu_B}{\sqrt{\sigma_B^2 + \varepsilon}} + \beta^{(l)}
\end{equation}
\begin{equation}
  \mathbf{a}^{(l)} = \text{ReLU}\!\left(\hat{\mathbf{z}}^{(l)}\right) = \max\!\left(0, \hat{\mathbf{z}}^{(l)}\right)
\end{equation}
\begin{equation}
  \mathbf{h}^{(l)} = \text{Dropout}_{p=0.30}\!\left(\mathbf{a}^{(l)}\right) = \mathbf{m}^{(l)} \odot \mathbf{a}^{(l)}
\end{equation}
where $\mathbf{m}^{(l)} \sim \text{Bernoulli}(1 - p)$ is a stochastic binary dropout mask.

The terminal output layer $L_5$ applies a linear affine projection without activation or
batch normalization to predict the continuous WQI value:
\begin{equation}
  \hat{y}_{\text{MLP}}(\mathbf{x}) = \mathbf{w}^{(5)\top}\mathbf{h}^{(4)} + b^{(5)}
\end{equation}

\subsection{Parameterization \& Weight Initialization}
\begin{itemize}
  \item \textbf{Total Trainable Parameters:} 207,681 weights and biases.
  \item \textbf{Weight Initialization:} Kaiming He normal initialization:
        \begin{equation}
          W_{ij}^{(l)} \sim \mathcal{N}\!\left(0, \; \sqrt{\frac{2}{d_{\text{in}}^{(l)}}}\right), \quad b_i^{(l)} = 0
        \end{equation}
        optimizing initial variance for ReLU non-linearities and preventing vanishing/exploding gradients.
  \item \textbf{Batch Normalization Running Momentum:} Set to $\alpha = 0.1$, $\varepsilon = 10^{-5}$.
\end{itemize}

\subsection{Secondary Network: ResidualMLPNet}
For the residual hybrid architecture (Task 4), a dedicated neural network,
\texttt{ResidualMLPNet}, was constructed to predict signed residual corrections.
It ingests 54 inputs (53 standardized DL features + 1 base tree prediction scalar)
with topology $54 \to 256 \to 128 \to 64 \to 1$ ($56,449$ parameters),
incorporating \texttt{BatchNorm1d} and Dropout ($p = 0.20$).

% ============================================================
\section{MLP Training, Optimization \& Uncertainty (Task 3)}
\label{sec:training}

\subsection{Optimization Configuration}
Training was executed on the full 1,973,730-record training partition:
\begin{itemize}
  \item \textbf{Optimizer:} AdamW~\cite{loshchilov2018decoupled} with decoupled weight decay ($\text{lr} = 10^{-3}$, $\lambda = 10^{-4}$, $\beta_1 = 0.9$, $\beta_2 = 0.999$).
  \item \textbf{Loss Function:} Mean Squared Error (MSE):
        \begin{equation}
          \mathcal{L}_{\text{MSE}} = \frac{1}{N}\sum_{i=1}^N \left(y_i - \hat{y}_i\right)^2
        \end{equation}
  \item \textbf{Batch Size:} 2,048 samples per mini-batch ($964$ iterations per epoch).
  \item \textbf{Epochs:} 3 full passes across 1.97M records ($2,892$ gradient updates).
  \item \textbf{Checkpointing:} Model weights were evaluated against the validation set
        ($N = 422,944$) at each epoch boundary; the checkpoint achieving minimum validation
        loss was restored for final test evaluation.
\end{itemize}

\subsection{Training Convergence Dynamics}
Fig.~\ref{fig:training_curves} illustrates the training loss and validation MAE convergence trajectories.
The model demonstrated rapid, stable optimization:
\begin{itemize}
  \item \textbf{Epoch 1:} Training loss dropped from initial $74.2$ to $3.82$; validation $\text{MAE} = 1.3412$.
  \item \textbf{Epoch 2:} Lowest validation loss reached ($2.6752$), validation $\text{MAE} = 1.2534$.
  \item \textbf{Epoch 3:} Validation $\text{MAE} = 1.2890$, confirming convergence without overfitting.
  \item \textbf{Held-Out Test Set Performance:} $\text{MAE} = 1.3800$, $\text{RMSE} = 1.6720$, $R^2 = 0.991029$.
\end{itemize}

\begin{figure}[!t]
  \centering
  \includegraphics[width=0.95\columnwidth]{figures/03_training_convergence_curves.png}
  \caption{Training loss and validation MAE convergence trajectories for \texttt{WaterQualityMLP} across training epochs, demonstrating stable optimization under Batch Normalization and Dropout ($p=0.30$).}
  \label{fig:training_curves}
\end{figure}

\subsection{Regularization \& Dropout Parameter Sweeps}
To evaluate regularization sensitivity, controlled experiments were conducted:
\begin{itemize}
  \item \textbf{Dropout Sweep ($p \in [0.0, 0.5]$):} As depicted in Fig.~\ref{fig:dropout_curves},
        $p = 0.0$ exhibited slight overfitting ($\text{Val MAE} = 1.521$), while $p = 0.50$
        caused underfitting ($\text{Val MAE} = 1.684$). The sweet spot emerged at $p = 0.30$
        ($\text{Val MAE} = 1.380$).
  \item \textbf{Weight Decay Sweep ($\lambda \in [0.0, 10^{-2}]$):} Weight decay $\lambda = 10^{-4}$
        provided optimal $L_2$ shrinkage. Larger penalties ($\lambda \ge 10^{-2}$) degraded
        predictive capacity by over-regularizing dense layer weights.
\end{itemize}

\begin{figure}[!t]
  \centering
  \includegraphics[width=0.95\columnwidth]{figures/07_regularization_dropout_curves.png}
  \caption{Validation MAE response across dropout rates ($p \in [0.0, 0.5]$), illustrating the regularization optimum at $p = 0.30$.}
  \label{fig:dropout_curves}
\end{figure}

\subsection{Monte Carlo Dropout Epistemic Uncertainty}
To provide calibrated predictive confidence, a Monte Carlo Dropout protocol was developed.
During test inference, $T = 30$ stochastic forward passes were executed with dropout active.

\noindent\textbf{Critical Implementation Fix:} In naive MC Dropout, executing \texttt{model.train()}
erroneously activates training mode in \texttt{BatchNorm1d} layers, causing batch statistics
to overwrite frozen running buffers ($\hat{\mu}_B, \hat{\sigma}_B^2$) and producing corrupted
predictions ($\text{MAE} = 3.1314$).
In our corrected implementation:
\begin{itemize}
  \item \texttt{model.eval()} is invoked globally, freezing all BatchNorm running statistics.
  \item Only \texttt{nn.Dropout} modules are explicitly switched to training mode:
        \texttt{for m in model.modules(): if isinstance(m, nn.Dropout): m.train()}.
\end{itemize}

\noindent\textbf{Corrected MC Dropout Benchmark} ($N = 1,000$ validation records, $T = 30$ stochastic passes):
\begin{itemize}
  \item \textbf{MC Ensemble Mean MAE:} $\mathbf{1.8079}$ (improved from $3.1314$ under corrupted BatchNorm).
  \item \textbf{MC Ensemble Mean RMSE:} $\mathbf{2.2994}$ (improved from $3.7004$).
  \item \textbf{MC Ensemble $R^2$:} $\mathbf{0.9825}$ (improved from $0.9547$).
  \item \textbf{Mean Epistemic Uncertainty ($\bar{\sigma}$):} $\mathbf{8.3666}$ WQI index units.
  \item \textbf{Median Epistemic Uncertainty ($\tilde{\sigma}$):} $\mathbf{8.4110}$ WQI index units.
  \item \textbf{Uncertainty Range:} $[\sigma_{\min} = 3.5040, \; \sigma_{\max} = 14.6220]$.
\end{itemize}

Observations exhibiting extreme geochemical excursions (dissolved oxygen deficit $> 6\text{ mg/L}$,
phosphorus spikes $> 1.5\text{ mg/L}$) generated epistemic uncertainty $\sigma > 12.0$,
functioning as an automatic reliability filter for anomalous environmental conditions.

% ============================================================
\section{Multi-Paradigm Residual Hybrid Architecture (Task 4)}
\label{sec:hybrid}

\subsection{Architectural Motivation}
Tabular environmental data with regulatory threshold boundaries presents an architectural dilemma:
decision trees partition continuous spaces along orthogonal step functions, perfectly matching
boolean thresholds in central data regions but producing step discretization errors on extreme tails.
Conversely, deep neural networks produce smooth continuous representations that reduce large
outlier errors but struggle with sharp threshold boundaries.

The \textbf{Residual Hybrid} architecture unites these complementary strengths:
\begin{equation}
  \hat{y}_{\text{hybrid}}(\mathbf{x}) = \hat{y}_{\text{base}}(\mathbf{x}) + \hat{r}(\mathbf{x})
  \label{eq:hybrid_formulation}
\end{equation}
where $\hat{y}_{\text{base}}$ is a frozen tree ensemble prediction and $\hat{r}(\mathbf{x})$ is a
neural network prediction of the signed residual error $y - \hat{y}_{\text{base}}$.

\subsection{Base Tree Ensemble Formulation}
The base prediction is computed as an equal-weighted convex combination of the Phase~1 models:
\begin{equation}
  \hat{y}_{\text{base}}(\mathbf{x}_{\text{tree}}) = 0.50\,\hat{y}_{\text{RF}}(\mathbf{x}_{\text{tree}}) + 0.50\,\hat{y}_{\text{XGB}}(\mathbf{x}_{\text{tree}})
  \label{eq:base_tree}
\end{equation}
This combination combines Random Forest's low-variance bagging with XGBoost's regularized gradient boosting.

\subsection{Leakage-Free 5-Fold OOF Training Protocol}
Training the residual network on in-sample base predictions ($\hat{y}_{\text{base}}(x_i)$ evaluated
on the training data) would cause catastrophic overfitting: the base trees have already memorized
training samples, producing near-zero residuals that do not reflect true test-time error distributions.

To prevent data contamination, we implemented a strict 5-fold Out-Of-Fold (OOF) protocol:
\begin{enumerate}
  \item The 1,973,730 training records were partitioned into 5 stratified folds.
  \item For each fold $k \in \{1,\dots,5\}$, base RF and XGBoost models were fit on folds $j \neq k$
        and evaluated on fold $k$.
  \item Unbiased out-of-fold residual targets were computed:
        \begin{equation}
          r_i^{\text{OOF}} = y_i - \hat{y}_{\text{base}}^{\text{OOF}}(\mathbf{x}_i) \quad \forall i \in \{1,\dots,N\}
        \end{equation}
  \item The secondary network, \texttt{ResidualMLPNet}, was trained to predict $r_i^{\text{OOF}}$
        from the 54-dimensional vector:
        \begin{equation}
          \mathbf{x}_{\text{res}} = \big[\,\mathbf{x}_{\text{DL}} \;\Vert\; \hat{y}_{\text{base}}^{\text{OOF}}\,\big] \in \mathbb{R}^{54}
        \end{equation}
        minimizing Smooth L1 (Huber) loss ($\beta = 1.0$) across 50 epochs:
        \begin{equation}
          \mathcal{L}_{\text{SL1}}(r, \hat{r}) =
          \begin{cases}
            0.5(r - \hat{r})^2 & \text{if } |r - \hat{r}| < 1.0 \\
            |r - \hat{r}| - 0.5 & \text{otherwise}
          \end{cases}
        \end{equation}
\end{enumerate}

The best validation checkpoint ($\text{Val MAE} = 0.0713$) was saved to
\texttt{models/hybrid\_ensemble\_checkpoint.pt}.

\begin{figure}[!t]
  \centering
  \includegraphics[width=0.95\columnwidth]{task_6_outputs/04_hybrid_component_contributions.png}
  \caption{Decomposition of hybrid predictions across test samples, showing base tree predictions ($\hat{y}_{\text{base}}$) providing the foundational estimate and the residual network providing subtle, bounded corrections ($\hat{r}$).}
  \label{fig:hybrid_components}
\end{figure}

\subsection{Empirical Residual Correction Mechanics}
Across the 422,943 test records, the neural residual corrections $\hat{r}$ exhibited:
\begin{itemize}
  \item \textbf{Range:} $[-7.3070, +3.2043]$ WQI index units.
  \item \textbf{Mean Correction:} $+0.0006$ units (unbiased zero-centering).
  \item \textbf{Standard Deviation:} $0.1611$ units.
  \item \textbf{Interquartile Range:} $[-0.0145, +0.0152]$ units.
\end{itemize}

As shown in Fig.~\ref{fig:correction_vs_error}, the neural residual correction $\hat{r}$
exhibits a strong positive empirical correlation ($r = 0.5262$) with the actual base tree error.
The residual network successfully reduces error in \textbf{50.67\%} of all test observations,
and in over \textbf{84.2\%} of high-error tail cases ($|\text{Error}_{\text{base}}| > 1.0$).

\begin{figure}[!t]
  \centering
  \includegraphics[width=0.95\columnwidth]{task_6_outputs/06_correction_vs_base_error.png}
  \caption{Empirical correlation ($r = 0.5262$) between base tree prediction error ($y - \hat{y}_{\text{base}}$) and neural residual correction ($\hat{r}$), proving targeted error correction on high-variance samples.}
  \label{fig:correction_vs_error}
\end{figure}

% ============================================================
\section{Comprehensive Model Evaluation \& Benchmarking (Task 5)}
\label{sec:results}

\subsection{Unified 6-Model Benchmark}
All six predictive architectures were evaluated on the common held-out test split
(`test.parquet`, $N = 422,943$ records). Table~\ref{tab:master_benchmark} presents
the empirical test set metrics.

\begin{table*}[!t]
  \caption{Comprehensive 6-Model Benchmark on Quarantined Held-Out Test Set ($N = 422,943$ Records)}
  \label{tab:master_benchmark}
  \centering
  \setlength{\tabcolsep}{4.5pt}
  \begin{tabular}{clcrrrrrrr}
    \toprule
    \textbf{Rank} & \textbf{Model Architecture} & \textbf{Paradigm} & \textbf{Input Dim}
      & \textbf{MAE} & \textbf{RMSE} & \textbf{$R^2$} & \textbf{MedAE} & \textbf{MaxAE} & \textbf{Latency} \\
    \midrule
    \textbf{1} & \textbf{Random Forest}      & Bagging Trees       & 38 Tabular          & \textbf{0.0514} & 0.2016          & 0.999870          & \textbf{0.0042} & 12.6128          & $18.8\text{ ms}$ \\
    \textbf{2} & \textbf{Residual Hybrid}    & Tree + Residual MLP & 38 Tab + 53 DL      & 0.0710          & \textbf{0.1678} & \textbf{0.999910} & 0.0229          & \textbf{5.6842}  & $24.7\text{ ms}$ \\
    \textbf{3} & \textbf{Base Tree Ensemble} & Equal Trees (50/50) & 38 Tabular          & 0.0900          & 0.2464          & 0.999805          & 0.0261          & 8.7128           & $22.8\text{ ms}$ \\
    \textbf{4} & \textbf{XGBoost}            & Gradient Boosting   & 38 Tabular          & 0.1614          & 0.4364          & 0.999389          & 0.0480          & 9.2769           & \textbf{4.6\,ms} \\
    \textbf{5} & \textbf{Weighted Ensemble}  & Convex Blend        & 38 Tab + 53 DL      & 0.3096          & 0.4155          & 0.999446          & 0.2653          & 8.8283           & $36.1\text{ ms}$ \\
    \textbf{6} & \textbf{WaterQualityMLP}    & Deep Neural Net     & 53 (36 num+17 OHE)  & 1.3800          & 1.6720          & 0.991029          & 1.2893          & 31.1327          & $14.3\text{ ms}$ \\
    \bottomrule
  \end{tabular}
\end{table*}

\begin{figure}[!t]
  \centering
  \includegraphics[width=0.95\columnwidth]{figures/05_model_performance_comparison.png}
  \caption{Unified 6-model benchmark on the held-out test set comparing MAE, RMSE, and $R^2$ across all evaluated architectures.}
  \label{fig:model_perf}
\end{figure}

\begin{figure}[!t]
  \centering
  \includegraphics[width=0.95\columnwidth]{task_6_outputs/09_actual_vs_predicted.png}
  \caption{Actual versus predicted CCME WQI scatter plots across the six model architectures on the 422,943-record test split.}
  \label{fig:actual_vs_pred}
\end{figure}

\subsection{Key Scientific Trade-Offs \& Findings}

\noindent\textbf{1.\ Peak Outlier Error Compression:}
While Random Forest achieves the lowest median error ($\text{MedAE} = 0.0042$), its
worst-case error reaches $\text{MaxAE} = 12.6128$ due to step-discretization errors.
The Residual Hybrid compresses the maximum absolute error down to $\mathbf{5.6842}$,
representing a \textbf{54.9\% reduction in peak prediction error}.

\noindent\textbf{2.\ Global Variance Minimization:}
Because RMSE heavily penalizes large deviation outliers, the Residual Hybrid achieves the
lowest RMSE across all evaluated architectures ($\mathbf{\text{RMSE} = 0.1678}$),
representing a \textbf{31.9\% variance reduction} over the Base Tree Ensemble ($0.2464$)
and achieving the highest coefficient of determination ($\mathbf{R^2 = 0.999910}$).

\noindent\textbf{3.\ Cumulative Error Distribution (CDF):}
Fig.~\ref{fig:error_cdf} plots the Cumulative Distribution Function of absolute errors.
Over $95\%$ of Residual Hybrid predictions exhibit absolute errors below $0.2039$ index units,
and $99\%$ fall below $0.6210$ units.

\begin{figure}[!t]
  \centering
  \includegraphics[width=0.95\columnwidth]{task_6_outputs/11_absolute_error_cdf.png}
  \caption{Cumulative Distribution Function (CDF) of absolute prediction errors across models on the held-out test split ($N = 422,943$).}
  \label{fig:error_cdf}
\end{figure}

\noindent\textbf{4.\ Tabular Inductive Bias Reality:}
The standalone MLP achieves $\text{MAE} = 1.3800$ and $R^2 = 0.991029$. While strong in
absolute terms, this represents a 27-fold higher MAE than Random Forest ($0.0514$).
The CCME WQI is calculated using sharp threshold step functions ($F_1, F_2, F_3$),
which axis-aligned decision trees approximate with fewer parameters than continuous neural networks.

\subsection{Pairwise Win-Rate Matrix}
Table~\ref{tab:winrate} reports the pairwise win-rate matrix, defined as the percentage
of test observations where the row model achieves lower absolute error than the column model.
Random Forest wins on typical median points ($61.4\%$ vs Hybrid), while the Residual Hybrid
wins decisively on high-error tail points ($|\text{error}| > 1.0$).

\begin{table}[!t]
  \caption{Pairwise Win-Rate Matrix (\% of Test Observations)}
  \label{tab:winrate}
  \centering
  \setlength{\tabcolsep}{3pt}
  \scriptsize
  \begin{tabular}{lrrrrrr}
    \toprule
    \textbf{Model} & \textbf{RF} & \textbf{Hybrid} & \textbf{Base Tree} & \textbf{XGB} & \textbf{Weighted} & \textbf{MLP} \\
    \midrule
    \textbf{Random Forest}     & 50.0\% & \textbf{61.4\%} & \textbf{78.2\%} & \textbf{90.8\%} & \textbf{96.9\%} & \textbf{97.7\%} \\
    \textbf{Residual Hybrid}   & 38.6\% & 50.0\%          & \textbf{68.4\%} & \textbf{84.2\%} & \textbf{94.5\%} & \textbf{96.8\%} \\
    \textbf{Base Tree Ens.}    & 21.8\% & 31.6\%          & 50.0\%          & \textbf{76.5\%} & \textbf{91.2\%} & \textbf{96.1\%} \\
    \textbf{XGBoost}           &  9.2\% & 15.8\%          & 23.5\%          & 50.0\%          & \textbf{81.4\%} & \textbf{93.5\%} \\
    \textbf{Weighted Ensemble} &  3.1\% &  5.5\%          &  8.8\%          & 18.6\%          & 50.0\%          & \textbf{92.1\%} \\
    \textbf{WaterQualityMLP}   &  2.3\% &  3.2\%          &  3.9\%          &  6.5\%          &  7.9\%          & 50.0\% \\
    \bottomrule
  \end{tabular}
\end{table}

% ============================================================
\section{Systematic 7-Group Feature Ablation Studies (Task 5)}
\label{sec:ablation}

\subsection{Protocol}
To isolate the marginal information contribution of each domain sub-system, a Leave-One-Group-Out
(LOGO) ablation protocol was conducted on $N = 50,000$ validation records. Each feature group
was neutralized by substituting its sample population mean, measuring the exact performance
degradation ($\Delta\text{MAE}$, $\Delta R^2$) relative to the full-feature baseline.

\begin{figure}[!t]
  \centering
  \includegraphics[width=0.95\columnwidth]{figures/06_feature_group_ablation_impact.png}
  \caption{Leave-One-Group-Out (LOGO) feature ablation impact ($\Delta\text{MAE}$) across feature groups for Random Forest, XGBoost, and WaterQualityMLP.}
  \label{fig:ablation_fig}
\end{figure}

\subsection{Ablation Results}
Table~\ref{tab:ablation} reports the empirical degradation metrics across all 7 official groups.

\begin{table}[!t]
  \caption{7-Group Leave-One-Group-Out (LOGO) Ablation Results}
  \label{tab:ablation}
  \centering
  \setlength{\tabcolsep}{2.5pt}
  \scriptsize
  \begin{tabular}{lrrrrr}
    \toprule
    \textbf{Ablated Feature Group} & \textbf{$k$}
      & \textbf{RF $\Delta$MAE} & \textbf{XGB $\Delta$MAE} & \textbf{MLP $\Delta$MAE} & \textbf{MLP $\Delta R^2$} \\
    \midrule
    \textit{Full Feature Baseline}  & 0 & 0.0000 & 0.0000 & 0.0000 & 0.0000 \\
    \textbf{\texttt{base\_sensors}}      & 8 & \textbf{+12.5749} & \textbf{+21.6473} & +0.4090 & -0.0097 \\
    \textbf{\texttt{log1p\_normalized}}  & 7 & +7.8952 & +1.6068 & \textbf{+12.1642} & \textbf{-0.8511} \\
    \textbf{\texttt{stoichiometry}}      & 7 & +0.0183 & +1.1525 & +0.0363 & -0.0007 \\
    \textbf{\texttt{env\_stress}}        & 1 & +0.0000 & +0.1952 & +0.5474 & -0.0101 \\
    \textbf{\texttt{contextual}}         & 2 & +0.0000 & +0.0009 & +0.4214 & -0.0103 \\
    \textbf{\texttt{history\_lags}}      & 6 & -0.0000 & +0.0406 & +0.0139 & -0.0002 \\
    \textbf{\texttt{temporal\_cyclical}} & 7 & -0.0001 & +0.0002 & +0.0299 & -0.0002 \\
    \bottomrule
  \end{tabular}
\end{table}

\subsection{Architectural Findings}
\begin{enumerate}
  \item \textbf{Raw Sensors are Irreplaceable for Trees:} Neutralizing \texttt{base\_sensors}
        triggers catastrophic failure in tree models, increasing XGBoost error by $+21.65$ MAE
        and collapsing $R^2$ into negative territory. Raw pH, DO, and phosphorus carry the
        primary predictive signal.
  \item \textbf{Log-Transform is Mandatory for Neural Networks:} Removing \texttt{log1p\_normalized}
        causes catastrophic failure for \texttt{WaterQualityMLP} ($\Delta\text{MAE} = +12.16$,
        $\Delta R^2 = -0.8511$). Because nutrient telemetry exhibits extreme skewness ($\gamma_1 > 40$),
        the $\log(1+x)$ transformation is an indispensable requirement for gradient descent stability.
  \item \textbf{Stoichiometric Synergy in Gradient Boosting:} XGBoost effectively leverages
        Benson-Krause DO deficit~\cite{benson1984concentration} and Redfield N:P ratio~\cite{redfield1934proportions}s, showing a $+1.15$ MAE degradation
        when stoichiometric features are ablated.
\end{enumerate}

% ============================================================
\section{Computational Profiling \& Latency (Task 5)}
\label{sec:compute}

Inference and training benchmarks were conducted on Apple Silicon ARM64 (macOS), Python 3.11,
PyTorch 2.14.1, CPU execution mode (Table~\ref{tab:compute} and Fig.~\ref{fig:comp_tradeoffs}).

\begin{table}[!t]
  \caption{Computational Benchmarks across Model Architectures}
  \label{tab:compute}
  \centering
  \setlength{\tabcolsep}{3pt}
  \scriptsize
  \begin{tabular}{lrrrr}
    \toprule
    \textbf{Performance Metric} & \textbf{Random Forest} & \textbf{XGBoost} & \textbf{WaterQualityMLP} & \textbf{Residual Hybrid} \\
    \midrule
    \textbf{Single Latency (B=1)} & 18.79\,ms & \textbf{4.60\,ms} & \textbf{0.30\,ms} & 24.72\,ms \\
    \textbf{Throughput (100k)}    & \textbf{876,262/s} & 205,824/s & 227,599/s & 138,864/s \\
    \textbf{Training Duration}    & 32.40\,s & \textbf{8.20\,s} & 58.60\,s (3 ep) & 59.80\,s (50 ep) \\
    \textbf{Model Size on Disk}   & 82.4\,MB & \textbf{1.8\,MB} & \textbf{0.8\,MB} & \textbf{1.2\,MB} \\
    \bottomrule
  \end{tabular}
\end{table}

\begin{figure}[!t]
  \centering
  \includegraphics[width=0.95\columnwidth]{figures/10_computational_tradeoffs.png}
  \caption{Inference latency versus test RMSE trade-offs across model architectures, illustrating XGBoost's efficiency advantage for edge deployment.}
  \label{fig:comp_tradeoffs}
\end{figure}

\noindent\textbf{MLP Mini-Batch Scaling Profile (100k records):}
Mini-batch throughput scales rapidly from $103,974$ samples/s (batch size 64) to plateau
near batch size $16,384$ at $\approx 358,000$ samples/s.

\noindent\textbf{Operational Deployment Recommendation:}
\begin{itemize}
  \item \textbf{Edge/IoT In-Situ Probes:} XGBoost ($4.60\text{ ms}$ latency, $1.8\text{ MB}$ footprint).
  \item \textbf{Data Center High-Throughput Batch:} Random Forest ($876,262$ records/s).
  \item \textbf{Mission-Critical Environmental Compliance:} Residual Hybrid (lowest RMSE $0.1678$,
        $54.9\%$ peak error compression).
\end{itemize}

% ============================================================
\section{Subgroup Stratification Analysis (Task 5)}
\label{sec:subgroup}

\subsection{Cross-National Geographic Generalization}
Models were evaluated across all 5 sovereign jurisdictions (Table~\ref{tab:geo} and Fig.~\ref{fig:geo_perf}).
Both Random Forest and Residual Hybrid achieve $R^2 > 0.9989$ in all nations. Ireland and USA
exhibit the lowest errors ($\text{MAE} < 0.02$) due to dense, standardized continuous automated probes.

\begin{table}[!t]
  \caption{Cross-National Geographic Subgroup Performance}
  \label{tab:geo}
  \centering
  \setlength{\tabcolsep}{2.5pt}
  \scriptsize
  \begin{tabular}{lrrrrr}
    \toprule
    \textbf{Jurisdiction} & \textbf{Test Samples} & \textbf{RF MAE} & \textbf{Hybrid MAE} & \textbf{XGB MAE} & \textbf{MLP MAE} \\
    \midrule
    \textbf{Ireland} & 35,342 & \textbf{0.0082} & 0.0197 & 0.0482 & 1.7125 \\
    \textbf{USA}     & 60,850 & \textbf{0.0078} & 0.0172 & 0.0849 & 1.7051 \\
    \textbf{China}   &  6,915 & \textbf{0.0184} & 0.0584 & 0.3785 & 1.9694 \\
    \textbf{England} & 319,239 & \textbf{0.0651} & 0.0871 & 0.1834 & 1.2673 \\
    \textbf{Canada}  &    597 & \textbf{0.1246} & 0.1431 & 0.3437 & 1.9613 \\
    \bottomrule
  \end{tabular}
\end{table}

\begin{figure}[!t]
  \centering
  \includegraphics[width=0.95\columnwidth]{figures/08_geographic_country_performance.png}
  \caption{Cross-national geographic performance comparison across sovereign jurisdictions, illustrating consistently high accuracy across all monitoring networks.}
  \label{fig:geo_perf}
\end{figure}

\subsection{Aquatic Habitat Stratification}
Across 12 aquatic habitat classifications (Fig.~\ref{fig:waterbody_perf}), Residual Hybrid
achieves $R^2 > 0.998$ in all strata. Lentic waters (Lakes: $\text{MAE} = 0.0381$) and lotic
transport (Rivers: $\text{MAE} = 0.0652$) exhibit highest precision. Sewage effluent
($\text{MAE} = 0.2145$) and agricultural drainage ditches ($\text{MAE} = 0.1045$) display
elevated variance due to episodic organic loading spikes.

\begin{figure}[!t]
  \centering
  \includegraphics[width=0.95\columnwidth]{figures/09_aquatic_waterbody_performance.png}
  \caption{Aquatic ecosystem habitat performance (Residual Hybrid) across 12 waterbody classifications, showing stable precision across lentic, lotic, and marine environments.}
  \label{fig:waterbody_perf}
\end{figure}

\subsection{Temporal Invariance}
Evaluation across historical decadal cohorts (Pre-1990: $\text{MAE} = 0.0741$; 1990--2005:
$\text{MAE} = 0.0632$; Post-2005: $\text{MAE} = 0.0738$) and seasonal cycles (Spring: $0.0712$;
Summer: $0.0694$; Autumn: $0.0721$; Winter: $0.0715$) proves strong temporal invariance
across 83 years of historical sensor transitions.

% ============================================================
\section{Axiomatic Explainability via Integrated Gradients (Task 6)}
\label{sec:explainability}

\subsection{Methodological Formulation}
Integrated Gradients computes feature attribution along the straight-line path between
a reference baseline $\mathbf{x}'$ and input $\mathbf{x}$:
\begin{equation}
  \text{IG}_i(\mathbf{x}) = (x_i - x_i') \times \int_{0}^{1} \frac{\partial F\!\left(\mathbf{x}' + \alpha(\mathbf{x} - \mathbf{x}')\right)}{\partial x_i} \, d\alpha
  \label{eq:ig_integral}
\end{equation}
Approximated via $m = 50$ Gauss-Legendre quadrature steps:
\begin{equation}
  \widehat{\text{IG}}_i(\mathbf{x}) = (x_i - x_i') \times \frac{1}{m}\sum_{k=1}^m \frac{\partial F\!\left(\mathbf{x}' + \frac{k}{m}(\mathbf{x} - \mathbf{x}')\right)}{\partial x_i}
\end{equation}

\subsection{Valid Modal Observed Reference Baseline}
An all-zero vector in one-hot categorical space represents an impossible real-world state.
To ensure chemical and mathematical validity, we established a \textbf{modal observed reference baseline}:
\begin{itemize}
  \item Numerical dimensions ($x'_{1..36}$): $0.0$ (mean of standardized training distribution).
  \item Categorical one-hot ($x'_{37..53}$): Active one-hot indicator for modal training classes
        (\textbf{England}, \texttt{country\_code=2}; \textbf{River}, \texttt{waterbody\_type\_code=8}).
  \item Baseline model prediction: $F(\mathbf{x}') = 86.7233$ WQI index units (clean, typical river water).
\end{itemize}

\subsection{Path Convergence \& Completeness Audit}
Vectorized autograd evaluation across $N = 1,000$ test observations ($50,000$ forward-backward passes)
executed in $0.46\text{ s}$ on CPU. The Completeness axiom was numerically verified:
\begin{itemize}
  \item \textbf{Mean Completeness Divergence ($\bar{\Delta}$):} $\mathbf{0.2529}$ WQI index units.
  \item \textbf{$P_{95}$ Completeness Divergence:} $\mathbf{1.2000}$ WQI index units.
  \item \textbf{Maximum Completeness Divergence:} $\mathbf{5.8780}$ WQI index units.
\end{itemize}
This establishes tight numerical convergence ($< 0.3\%$ relative error on the CCME $[0, 100]$ scale).

\begin{figure}[!t]
  \centering
  \includegraphics[width=0.95\columnwidth]{task_6_outputs/01_mlp_global_feature_importance.png}
  \caption{Global feature importance for \texttt{WaterQualityMLP} computed via Integrated Gradients ($N=1,000$, $m=50$), aggregated across 38 domain features.}
  \label{fig:ig_global}
\end{figure}

\begin{figure}[!t]
  \centering
  \includegraphics[width=0.95\columnwidth]{task_6_outputs/02_mlp_signed_feature_attributions.png}
  \caption{Signed Integrated Gradients feature attributions, confirming physical realism: chemical contaminants negatively impact WQI, while dissolved oxygen compliance elevates it.}
  \label{fig:ig_signed}
\end{figure}

\subsection{Global Feature Rankings \& Directional Chemistry}
As shown in Table~\ref{tab:ig_rankings}, Fig.~\ref{fig:ig_global}, and Fig.~\ref{fig:ig_signed},
the top predictive drivers are:
\begin{enumerate}
  \item \texttt{log1p\_orthophosphate} ($|\text{IG}| = 5.8602$, signed $\text{IG} = -1.3428$):
        Dominant nutrient pollutant; elevated phosphate strongly suppresses predicted WQI.
  \item \texttt{log1p\_nitrogen} ($|\text{IG}| = 3.4170$, signed $\text{IG} = -0.4215$):
        Key agricultural runoff indicator.
  \item \texttt{log1p\_np\_ratio} ($|\text{IG}| = 3.0399$, signed $\text{IG} = -0.8541$):
        Stoichiometric eutrophication limitation factor.
  \item \texttt{log1p\_ammonia} ($|\text{IG}| = 1.7096$, signed $\text{IG} = -1.1318$):
        Direct biological toxicant.
  \item \texttt{exceedance\_counter} ($|\text{IG}| = 0.6247$, signed $\text{IG} = -0.2449$):
        Regulatory violation tracking.
\end{enumerate}

\begin{table}[!t]
  \caption{Top 10 Global Integrated Gradients Feature Attributions}
  \label{tab:ig_rankings}
  \centering
  \setlength{\tabcolsep}{3pt}
  \scriptsize
  \begin{tabular}{lrrl}
    \toprule
    \textbf{Feature Dimension (53D)} & \textbf{Mean $|\text{IG}|$} & \textbf{Mean Signed $\text{IG}$} & \textbf{Environmental Direction} \\
    \midrule
    \texttt{log1p\_orthophosphate} & \textbf{5.8602} & $-1.3428$ & Pollutant $\downarrow$ WQI \\
    \texttt{log1p\_nitrogen}       & 3.4170 & $-0.4215$ & Pollutant $\downarrow$ WQI \\
    \texttt{log1p\_np\_ratio}      & 3.0399 & $-0.8541$ & Eutrophication $\downarrow$ WQI \\
    \texttt{log1p\_ammonia}        & 1.7096 & $-1.1318$ & Toxicant $\downarrow$ WQI \\
    \texttt{exceedance\_counter}   & 0.6247 & $-0.2449$ & Violations $\downarrow$ WQI \\
    \texttt{log1p\_bod}            & 0.5412 & $-0.5264$ & Organic load $\downarrow$ WQI \\
    \texttt{bod\_do\_ratio}        & 0.4812 & $-0.3120$ & Stress ratio $\downarrow$ WQI \\
    \texttt{do\_deficit}           & 0.4510 & $-0.2915$ & Hypoxia $\downarrow$ WQI \\
    \texttt{country\_code\_2} (England) & 0.3120 & $+0.1045$ & Context $\uparrow$ WQI \\
    \texttt{waterbody\_type\_8} (River) & 0.2415 & $+0.0812$ & Context $\uparrow$ WQI \\
    \bottomrule
  \end{tabular}
\end{table}

\subsection{Individual Prediction Explanations}
Fig.~\ref{fig:ig_individual} illustrates local additive feature decompositions for representative test observations:
\begin{itemize}
  \item \textbf{Excellent Water Quality} (Actual: $98.40$, Pred: $98.15$, Residual: $-0.25$): Low orthophosphate ($+4.82$) and zero guideline exceedances ($+1.21$) produce strong positive attributions.
  \item \textbf{Poor Water Quality} (Actual: $38.20$, Pred: $37.90$, Residual: $-0.30$): Orthophosphate spikes ($-14.85$) and ammonia toxicity ($-6.42$) push predicted WQI into the severely contaminated regime.
\end{itemize}

\begin{figure}[!t]
  \centering
  \includegraphics[width=0.95\columnwidth]{task_6_outputs/03_individual_prediction_explanations.png}
  \caption{Local Integrated Gradients feature attribution waterfall plots for representative test observations (Excellent vs Poor water quality).}
  \label{fig:ig_individual}
\end{figure}

% ============================================================
\section{Prediction Error Analysis \& Residual Diagnostics (Task 6)}
\label{sec:errors}

\subsection{Standard Residual Sign Convention}
The standard engineering residual convention is strictly enforced:
\begin{equation}
  \text{Residual} = \hat{y} - y = \text{Prediction} - \text{Actual}
\end{equation}
Positive residuals indicate model overprediction ($\hat{y} > y$); negative residuals indicate underprediction ($\hat{y} < y$).

\subsection{Residual Distribution Diagnostics}
As shown in Fig.~\ref{fig:res_dist} and Fig.~\ref{fig:res_box}, tree models and the Residual Hybrid
exhibit zero-centered bias ($\mu_e \in [+0.0002, +0.0011]$). The Residual Hybrid significantly
compresses the outlier whiskers relative to pure tree models.

\begin{figure}[!t]
  \centering
  \includegraphics[width=0.95\columnwidth]{task_6_outputs/07_residual_distributions_comparison.png}
  \caption{Kernel density estimation of prediction residual distributions ($\hat{y} - y$) on 422,943 test records, illustrating zero-centered bias across models.}
  \label{fig:res_dist}
\end{figure}

\begin{figure}[!t]
  \centering
  \includegraphics[width=0.95\columnwidth]{task_6_outputs/08_residual_boxplots.png}
  \caption{Residual error boxplots illustrating extreme outlier error compression achieved by the Residual Hybrid model.}
  \label{fig:res_box}
\end{figure}

\subsection{Error Stratification across CCME WQI Quality Tiers}
As shown in Fig.~\ref{fig:error_tier}, prediction errors are concentrated in the \textbf{Poor}
quality tier ($WQI \in [0, 45]$), where chemical parameters reach extreme, non-linear concentrations:
\begin{itemize}
  \item \textbf{Excellent Tier} ($N=153,362$): Hybrid $\text{MAE} = 0.0691$, Random Forest $\text{MAE} = 0.0502$.
  \item \textbf{Good Tier} ($N=169,842$): Hybrid $\text{MAE} = 0.0684$, Random Forest $\text{MAE} = 0.0489$.
  \item \textbf{Fair Tier} ($N=85,412$): Hybrid $\text{MAE} = 0.0721$, Random Forest $\text{MAE} = 0.0541$.
  \item \textbf{Marginal Tier} ($N=12,845$): Hybrid $\text{MAE} = 0.0784$, Random Forest $\text{MAE} = 0.0624$.
  \item \textbf{Poor Tier} ($N=1,482$): Hybrid $\text{MAE} = 0.0915$, Random Forest $\text{MAE} = 0.0812$, MLP $\text{MAE} = 1.8412$.
\end{itemize}

\begin{figure}[!t]
  \centering
  \includegraphics[width=0.95\columnwidth]{task_6_outputs/12_error_by_wqi_tier.png}
  \caption{Mean Absolute Error stratified across CCME water quality classification tiers, demonstrating stability across all tiers with higher variance in severely degraded waters.}
  \label{fig:error_tier}
\end{figure}

\subsection{Representative Difficult Cases}
\noindent\textbf{1.\ Extreme Hybrid Error Recovery Cases:}
\begin{itemize}
  \item \textbf{Case A (River, England):} Actual WQI $= 72.40$. Base Tree Predicted $= 81.11$ ($\text{Error} = 8.71$). Residual Hybrid Predicted $= 74.22$ ($\text{Error} = \mathbf{1.82}$, a $+6.89$ unit correction).
  \item \textbf{Case B (Effluent, England):} Actual WQI $= 64.80$. Base Tree Predicted $= 71.95$ ($\text{Error} = 7.15$). Residual Hybrid Predicted $= 66.12$ ($\text{Error} = \mathbf{1.32}$, a $+5.83$ unit correction).
\end{itemize}

\noindent\textbf{2.\ Extreme MLP Failures:}
\begin{itemize}
  \item \textbf{Overprediction Case:} Actual WQI $= 61.20$, Predicted $= 88.42$, Residual $= \mathbf{+27.22}$. Normal DO masked heavy organic contaminant loading.
  \item \textbf{Underprediction Case:} Actual WQI $= 94.50$, Predicted $= 63.37$, Residual $= \mathbf{-31.13}$. High localized orthophosphate triggered excessive continuous penalty.
\end{itemize}

% ============================================================
\section{Discussion \& Scientific Limitations}
\label{sec:limitations}

\subsection{Deterministic Target Approximation}
A critical qualification must accompany the near-perfect regression scores ($R^2 > 0.999$):
The CCME WQI is calculated through a deterministic mathematical formula based on discrete guideline exceedances.
Consequently, tree models act as efficient piecewise function approximators. These high coefficients
of determination demonstrate accurate reconstruction of regulatory policy formulas rather than
unconstrained natural process forecasting across unmonitored watersheds.

\subsection{Feature Leakage Considerations}
The \texttt{exceedance\_counter} feature counts guideline violations, mirroring the $F_1$ scope
component of the CCME formula. However, the LOGO ablation study (Section~\ref{sec:ablation})
definitively proves that when \texttt{exceedance\_counter} is completely removed, tree models
retain $R^2 > 0.999$, demonstrating that predictive performance is anchored in raw physical telemetry.

\subsection{Identified Scientific Limitations}
\begin{enumerate}
  \item \textbf{Geographic Representation Imbalance:} England provides $83.74\%$ of all observations,
        while Canada accounts for only $0.20\%$ ($597$ test records). Cross-national findings for
        Canada should be interpreted with appropriate caution.
  \item \textbf{Hybrid Inference Latency:} The Residual Hybrid's sequential inference pipeline
        ($24.72\text{ ms}$) involves executing two tree models plus the PyTorch neural network,
        making it 5$\times$ slower than standalone XGBoost ($4.60\text{ ms}$).
  \item \textbf{Tabular Deep Learning Gap:} Standalone \texttt{WaterQualityMLP} exhibited an MAE of
        $1.3800$ --- 27$\times$ higher than Random Forest ($0.0514$) --- highlighting the inherent
        difficulty of gradient-based optimization on discontinuous tabular threshold boundaries.
\end{enumerate}

% ============================================================
\section{Phase 2 Software Architecture \& Testing}
\label{sec:software}

Phase~2 delivers a modular Python library under \texttt{Phase\_2\_Deep\_Learning/src/}:
\begin{itemize}
  \item \texttt{src/dataset.py} --- \texttt{WaterQualityDataset} class wrapping tensored DL features.
  \item \texttt{src/dl\_models.py} --- \texttt{WaterQualityMLP} and \texttt{ResidualMLPNet} PyTorch definitions.
  \item \texttt{src/training.py} --- Training loops, early stopping, AdamW optimization, checkpointing.
  \item \texttt{src/hybrid\_models.py} --- \texttt{ResidualHybridEnsemble} inference pipeline.
  \item \texttt{src/oof\_generator.py} --- 5-fold stratified OOF residual generator.
  \item \texttt{src/evaluation.py} --- Metric calculators, subgroup stratification, win-rate matrix.
\end{itemize}

An automated test suite in \texttt{tests/} executes \textbf{104 unit tests} under \texttt{pytest},
covering model architecture forward passes, dataset tensor shapes, hybrid inference arithmetic,
and zero-leakage quarantine verification with \textbf{100\% passing status}.

% ============================================================
\section{Team Contribution Matrix \& Viva Voce Defense}
\label{sec:team}

In compliance with CSAI project assessment guidelines, Table~\ref{tab:team_contributions}
defines the Phase~2 work distribution and viva voce defense specialization across all six team members.

\begin{table}[!t]
  \caption{Phase 2 Technical Work Packages \& Viva Defense Specialization Matrix}
  \label{tab:team_contributions}
  \centering
  \setlength{\tabcolsep}{2.5pt}
  \scriptsize
  \begin{tabular}{p{1.5cm}p{2.6cm}p{3.9cm}}
    \toprule
    \textbf{Member} & \textbf{Phase 2 Technical Work Packages} & \textbf{Viva Defense Specialization} \\
    \midrule
    \textbf{Aman Soni}
      & \textbf{Tasks 1 \& 6: DL Pipeline \& Explainability}
      & PyTorch DataLoaders, 38 vs 53 feature representations, Integrated Gradients path integration, Completeness axiom, residual error diagnostics. \\
    \midrule
    \textbf{Mayank Sharma}
      & \textbf{Task 2: MLP Architecture Design}
      & Neural network topology ($53\to 512\to 256\to 128\to 64\to 1$), Kaiming He initialization, BatchNorm1d vs Dropout interaction, tabular inductive bias. \\
    \midrule
    \textbf{Chaitanya Kumar}
      & \textbf{Task 3: MLP Training \& Optimization}
      & AdamW optimizer, MSE vs Smooth L1 loss, validation loss curves, MC Dropout BatchNorm fix, epistemic uncertainty quantification. \\
    \midrule
    \textbf{Ankit Kumar Pandey}
      & \textbf{Task 4: Residual Hybrid Development}
      & Base tree ensembling ($0.5\text{RF}+0.5\text{XGB}$), 5-fold OOF residual generation, ResidualMLPNet (54-dim input), $54.9\%$ peak error compression. \\
    \midrule
    \textbf{Aaryan Gera}
      & \textbf{Task 5: Model Evaluation \& Ablation}
      & 6-model benchmark, 7-group LOGO ablation protocol, latency profiling (CPU/GPU), geographic and aquatic habitat subgroup stratification. \\
    \midrule
    \textbf{Pranav Sehgal}
      & Technical Documentation \& Review
      & CCME index mathematical formulation, environmental chemistry references, scientific limitations, Phase 2 technical report coordination. \\
    \bottomrule
  \end{tabular}
\end{table}

% ============================================================
\section{Conclusion}
\label{sec:conclusion}

Phase~2 delivered three interconnected deep learning contributions to continuous water quality prediction.
The \texttt{WaterQualityMLP} demonstrated that feed-forward neural networks achieve strong general approximation
($R^2 = 0.991$) on structured environmental telemetry, though an architectural gap persists relative to tree
ensembles on tabular step-function thresholds. The novel \textbf{Residual Hybrid} architecture successfully
bridged this gap: by uniting the orthogonal threshold precision of tree ensembles with the smooth manifold
error correction of a deep residual network, it achieved the \textbf{lowest global RMSE ($0.1678$)} across all
six evaluated paradigms and compressed worst-case prediction errors by \textbf{54.9\%} (MaxAE from $12.61$ to $5.68$).

Systematic 7-group LOGO ablation established that raw physical sensors are irreplaceable for tree models,
while logarithmic variance stabilization is mathematically essential for neural network gradient stability.
Axiomatic Integrated Gradients confirmed complete physical realism, demonstrating that orthophosphate,
nitrogen, and stoichiometric $\text{N}:\text{P}$ ratios serve as the dominant negative drivers of aquatic health.
Monte Carlo Dropout epistemic uncertainty quantification enabled automated out-of-distribution sample detection
for extreme geochemical excursions.

Collectively, Phase~2 establishes that multi-paradigm residual hybridization represents the optimal framework
for real-world environmental monitoring systems, providing unprecedented precision, robustness against outlier
errors, and rigorous axiomatic interpretability.

% ── References ────────────────────────────────────────────────
\balance
\bibliographystyle{IEEEtran}
\bibliography{references_p2}

\end{document}
